%%%%%%%%%%%%%%%%%%%%%%%%%%%%%%%%%%%%%%%%%%%%%%%%%%%%%%%%%%%%%%%%%%%%%%%%%%%
\chapter{Results and Discussion}\label{chap:results}
%%%%%%%%%%%%%%%%%%%%%%%%%%%%%%%%%%%%%%%%%%%%%%%%%%%%%%%%%%%%%%%%%%%%%%%%%%%

\section{Normal DQN Performance Ceiling}

The multi-phase curriculum concluded with 37,213 episodes of self-play training. The normal DQN agent used an LSTM-based history encoder and trained against random, greedy heuristic, league, and self-play opponents. Table~\ref{tab:vanilla_performance} reports the results.

\begin{table}[htbp]
    \centering
    \caption{Normal DQN performance after 37,213 episodes of self-play training.}
    \label{tab:vanilla_performance}
    \begin{tabular}{|l|c|c|c|c|}
        \hline
        \textbf{Opponent Type} & \textbf{Episodes} & \textbf{Wins} & \textbf{Win Rate} & \textbf{Draw Rate} \\
        \hline
        Greedy Heuristic & 9,319 & 1,418 & 15.2\% & 41.3\% \\
        Self-Play        & 9,565 & 2,370 & 24.8\% & 36.8\% \\
        Random           & 9,294 & 2,260 & 24.3\% & 37.7\% \\
        League           & 9,035 & 2,231 & 24.7\% & 38.2\% \\
        \hline
        \textbf{Aggregate} & \textbf{37,213} & \textbf{8,279} & \textbf{22.2\%} & \textbf{38.5\%} \\
        \hline
    \end{tabular}
\end{table}

The training data contains six findings, each linked to an architectural limitation.

\subsection{Win Rate Plateau}\label{finding:1}
Agent1's cumulative win rate stayed near 22\% after 2,000 episodes, it stayed more or less the same over the next 35,000. The 100-episode rolling average swung between 15\% and 30\% with no deviation.

\begin{figure}[htbp]
    \centering
    \includegraphics[width=0.85\textwidth]{figures/Results/cropped/cumulative_wins_winrate.png}
    \caption{Cumulative win count and win rate over 37,213 episodes. The win rate plateaus early and remains flat, indicating policy stagnation.}
    \label{fig:cumulative_winrate}
\end{figure}

\subsection{Loss--Performance Decoupling}\label{finding:2}
The TD error fell to roughly 0.001 however, the win rate did not follow. The network learned to predict the value of passive, draw-seeking play accurately. It did not discover better strategies.

\begin{figure}[htbp]
    \centering
    \includegraphics[width=0.85\textwidth]{figures/Results/cropped/policy_loss.png}
    \caption{Policy loss (TD error).}
    \label{fig:policy_loss}
\end{figure}

\subsection{Q-Value Overestimation}\label{finding:3}
Average Q-values climbed from 0.2 to 1.2. Win rates did not move with the rising value estimates.

\begin{figure}[htbp]
    \centering
    \includegraphics[width=0.85\textwidth]{figures/Results/cropped/q_value.png}
    \caption{Average Q-values climb.}
    \label{fig:q_value}
\end{figure}

\subsection{LSTM Prediction Accuracy Ceiling}\label{finding:4}
The LSTM piece-identity module accuracy stayed around 18--22\%, a 2.2--2.6$\times$ gain over the 8.3\% random baseline but is still far short of reliable tactical inference. The improvement stopped after episode 5,000 and stayed at that level.

\begin{figure}[htbp]
    \centering
    \includegraphics[width=0.85\textwidth]{figures/Results/cropped/prediction_accuracy.png}
    \caption{Piece-identity prediction accuracy.}
    \label{fig:prediction_accuracy}
\end{figure}

\subsection{Opponent-Specific Performance Asymmetry}\label{finding:5}
Performance varies by opponent given randomly in the training process. Win rate against the greedy heuristic was 15.2\%; in self-play it reached 24.8\%. The policy fit its own behavioral distribution and did not adapt to the others opponent types.

\subsection{Draw Dominance}\label{finding:6}
The game results was mostly draws. Draws accounted for 38.5\% of outcomes, far above the 22.2\% win rate. The agent learned to shuffle pieces passively and avoid uncertain engagements. It minimized short-term losses at the expense of long-term winning probability, lacking the long horizon planning.




\section{Behavioral Patterns from Training Logs}

\subsection{Passive Shuffling Convergence}

The agent grew conservative. Average episode length rose throughout training as it avoided combat and cycled pieces between adjacent squares. Scalar value collapse (Finding~2) explains why the Q-network converges to the lowest-variance action class, playing safe moves instead of winning moves.

\subsection{Opponent-Specific Adaptation Failure}

Table~\ref{tab:vanilla_performance} shows the agent overfitting to its own play style, not generalizing to other agents, since the drop-off of performance against greedy opponents is significant (15.2\%). The greedy heuristic has the lowest win rate (15.2\%) due to its rank-agnostic attack pattern directly punishing passive shuffling. The heuristic set of instructions attacks the nearest piece, forcing repeated engagements that the agent's risk-averse policy is designed to avoid. When forced into combat, the agent unreliable piece-identity inference needed to choose favorable matchups fails, so each engagement is approximately a coin flip weighted by the opponent's rank distribution. The result is a consistent loss rate that the agent win with further passivity.

Self-play produced the highest win rate (24.8\%), but this figure can be misleading. Both agents share the same weights and therefore the same behavioral tendencies. When two equally passive and cautious policies face each other, games resolve through symmetric attrition rather than strategic superiority. The apparent 24.8\% win rate reflects the probability that stochastic tiebreakers favor one copy over another, not genuine tactical competence. The high draw rate in self-play (36.8\%) confirms that most games end in mutual passivity.

Random opponents produced a 24.3\% win rate, nearly identical to self-play despite being the weakest opponent class. A random policy does not exploit the agent's weaknesses, but it also provides no consistent pattern for the agent to learn from. Transitions generated against random play carry high variance and low signal, contributing gradient noise without driving policy improvement. The similar win rates against random and self-play opponents indicate that the agent has not learned to exploit weak play, further evidence of policy stagnation.

League opponents (24.7\%) are a historical snapshots of the agent past versions. Because the agent's policy plateaued early (Finding~1), these snapshots are nearly indistinguishable from the current policy. A league of stagnant policies cannot force adaptation and the fixed $\epsilon$-greedy decay adds to this problem by providing no mechanism for adapting exploration to different opponent since the agent has stopped trying to find different strategy.

\subsection{Loss-Reward Divergence}

The TD loss dropped from 0.05 to 0.001 by episode 37,000. Average episode reward stayed flat near 0.08, a possible confusion in the reward signal. The network learned to predict outcomes of its existing passive play with high precision, knowing that a safe move will net a safe outcome. With the bottoming $\epsilon$-greedy, it did not search for strategies with higher expected returns.

\subsection{LSTM Accuracy}

The piece prediction accuracy hit 20\% after episode 5,000 and stagnated for the rest of the training. One correct inference in five is not enough to evaluate combat outcomes reliably, so the agent tend to avoid combat altogether


\section{Draw Dominance}

Draws reached 38.5\% of all games, and the win rate was only 22.2\%. This rate matters not just as a performance number but as a drag on learning itself, a symptom of a bigger problem.

The scalar value collapse (Finding~2) made the model play safe rather than optimal. A scalar Q-value cannot represent multi-modal outcome distributions. A move with 40\% win probability and 60\% loss probability (expected value $\approx -0.2$) looks worse than a certain draw ($-0.3$) to a mean-based estimator. The draw-producing move gets the higher Q-value even though it is strategically inferior.

\section{Architectural Comparison: Unified versus Dual-Model Pipeline}

A different approach was use from the initial architecture, which was using two separate networks, a RLNN SetupAgent for piece placement and a DQN for movement. That dual-model design was abandoned, due to the computation overhead and signal interference between the two networks. The current unified pipeline uses a single DQN for all movement and a randomized heuristic for placement. Table~\ref{tab:architecture_comparison} compares both.

\begin{table}[htbp]
    \centering
    \caption{Dual-model versus unified architectural pipelines.}
    \label{tab:architecture_comparison}
    \begin{tabular}{|p{3cm}|p{5.5cm}|p{5.5cm}|}
        \hline
        \textbf{Aspect} & \textbf{Dual-Model (Setup + Attack)} & \textbf{Unified DQN} \\
        \hline
        Architecture & Separate C51-based SetupAgent for placement, separate DQN for attack & Single DQN for movement; heuristic setup with randomized flag \\
        \hline
        Setup Latency & Approximately 2 seconds per episode & Under 100 milliseconds per episode \\
        \hline
        Training Overhead & Doubled compute requirement (setup and attack training phases) & Halved total training compute \\
        \hline
        Exploitability & Deterministic setup patterns learned by opponents & Randomized placement resists opponent exploitation \\
        \hline
        Draw Pattern & Draws due to gradient interference between competing objectives & Draws due to risk-averse policy under uncertainty \\
        \hline
    \end{tabular}
\end{table}

Two lessons emerged from this transition. The SetupAgent did not converged on deterministic placements, randomizing setup while lowering the TD-error of the model a mismatch in result. Placing flag in front of the lineup and trapping big pieces behind the bomb. A randomized heuristic eliminated this vulnerability and cut setup latency by 20$\times$.

The dual-model design also suffered from gradient interference. Setup gradients rewarded defensive formations while attack gradients rewarded aggressive piece use. Optimal defense often places pieces in positions bad for attack. Removing the setup phase from the learning loop resolved this conflict.

Both architectures converged on draw-heavy outcomes. In the dual case, gradient interference between objectives prevented coherent learning. In the unified case, scalar value collapse and poor exploration drove the agent into passive shuffling. The draw problem is architectural and rooted in the long horizon of the game, not an unforeseen accident.


\section{Architectural Bottleneck Analysis}

Each of the findings above can be explained by a known structural property of the DQN architecture. The explanations are hypotheses that the later sections test; they do not rule out other causes, such as the reward design (Appendix~\ref{app:lessons}).

\subsection{Q-Value Overestimation}

Finding~3 is a known property of single-estimator Q-learning. The network selects and evaluates actions with the same weights:

\begin{equation}
y_t = r_t + \gamma \max_{a'} Q_\theta(s_{t+1}, a')
\end{equation}

Noisy estimates plus a max operator produce a positive bias, reducing the ability of the agent to learn from its mistake \cite{RainbowJMLR2021}:

\begin{equation}
\mathbb{E}\left[\max_{a} (Q^*(s, a) + \epsilon_a)\right] \geq \max_{a} Q^*(s, a)
\end{equation}

The bias grows with noise variance and action-space size $|\mathcal{A}|$. Stratego positions can have over in theory a 220 to 240 legal moves, amplifying the effect. Double DQN splits selection from evaluation across two estimators, provably reducing this bias \cite{RainbowJMLR2021}.

\subsection{Scalar Value Collapse}

Finding~2 (loss convergence without performance gain) follows from distributional RL theory. The standard DQN learns a scalar $\mathbb{E}[Z(s,a)]$ and minimizes:

\begin{equation}
\mathcal{L} = \mathbb{E}\left[\left(r + \gamma \max_{a'} Q(s', a') - Q(s, a)\right)^2\right]
\end{equation}

In Stratego the return distribution $Z(s,a)$ is multi-modal \cite{RainbowJMLR2021}: wins near $+1$, losses near $-1$, exchanges in $[-0.5, +0.5]$, draws near $0$. A scalar mean collapses all modes into one number that may not match any real outcome. The loss converges when this mean is estimated correctly, even if the policy cannot tell aggressive high-variance moves from passive low-variance ones.


\subsection{LSTM Information Bottleneck}

Finding~4 follows from sequential compression. The LSTM encodes observations $o_1, \ldots, o_T$ through a 64-dimensional hidden state:

\begin{equation}
h_t = \text{LSTM}(o_t, h_{t-1})
\end{equation}

A Stratego game runs 200 to 500 moves. The hidden state must track belief estimates for 40 opponent pieces across 12 types. A 64-dimensional vector can hold enough information in theory, but every successive nonlinear update risks losing early observations. Information from move $k$ must survive $T - k$ transformations to reach the final representation. Gating helps, but retention still degrades with sequence length \cite{Aaren}.

AAREN sidesteps this bottleneck. It computes weighted aggregations directly over the full history with $O(1)$ per-timestep complexity:

\begin{equation}
h_t = \frac{a_t}{c_t} = \frac{\sum_{k=1}^{t} v_k \exp(s_k - m_t)}{\sum_{k=1}^{t} \exp(s_k - m_t)}
\end{equation}

No sequential compression, no information decay. The network can attend to any prior timestep directly.

\subsection{Exploration Collapse}

The normal DQN relies on a hard-coded $\epsilon$-greedy exploration schedule that decays from 1.0 down to 0.01. Once it hits this floor, the schedule stops adapting. From that point on, the agent explores uniformly at a flat 1\% rate regardless of the game state or its confidence in the current Q-values.

Stratego has over $10^{535}$ possible states \cite{perolat2022stratego}. A uniform random exploration rate of 1\% is a blunt mechanism that will almost never stumble on the complex sequences of moves required to discover optimal strategies. The agent becomes trapped in local optima because its exploration policy cannot recognize which parts of the state space require further investigation. This limitation necessitates a mechanism like Noisy Networks \cite{RainbowJMLR2021}, which replaces hard-coded randomness with learnable parameters that concentrate exploration where the value function is most uncertain.

\subsection{Sparse Reward Propagation}

The normal DQN cannot develop long-horizon strategic behavior (Findings~1, 5, 6), a limitation compounded by single-step TD learning. Terminal reward information at step $T$ propagates backward through the Bellman backup at a rate set by the discount factor $\gamma$:

\begin{equation}
Q(s_k, a_k) \leftarrow Q(s_k, a_k) + \alpha \left( r_k + \gamma Q(s_{k+1}, a_{k+1}) - Q(s_k, a_k) \right)
\end{equation}

For a decisive game outcome occurring at step $T$ in a game of length 300, a mid-game decision at step 100 receives an effective signal of $\gamma^{200}$. With $\gamma = 0.99$, this yields $0.99^{200} \approx 0.134$, meaning that only 13.4\% of the terminal reward signal reaches mid-game decisions through single-step bootstrapping. Multi-step returns ($n$-step TD) \cite{RainbowJMLR2021} propagate terminal signals $n$ times faster by directly incorporating rewards over $n$ consecutive steps:

\begin{equation}
y_t^{(n)} = \sum_{k=0}^{n-1} \gamma^k r_{t+k} + \gamma^n Q(s_{t+n}, a_{t+n})
\end{equation}

This reduces the effective bootstrap depth and accelerates the propagation of sparse terminal rewards to earlier decision points.


\section{Empirical Validation: Rainbow DQN with AAREN}

The bottleneck analysis above predicted that each failure mode has a specific architectural fix. The following data mining analysis tests those predictions empirically. Three training datasets were analyzed: the same normal DQN at 15,302 episodes, the Rainbow DQN with AAREN at 9,778 episodes of self-play, and a separate extended Rainbow run of 75,047 episodes against heuristic opponents. Table~\ref{tab:arch_comparison} presents the head-to-head comparison. The complete training progress charts appear in Appendix B (Figures~\ref{fig:lstm15k_progress}--\ref{fig:rdqn75k_aaren}).

Three limits apply to every comparison in this section. First, the runs differ in more than the architecture: the opponents, curriculum, reward and run length differ too. Each configuration was also trained once, so a difference cannot be attributed to a single component. Second, the self-play figures come from the training games, in which the agent plays a copy of itself, not from evaluations against fixed opponents. Third, the average Q-value logged by the Rainbow trainer was the maximum over all 3,600 actions. Most of those actions are illegal and never trained, and Appendix~\ref{app:lessons} shows this statistic read about $+2$ while the returns were about $-1$. The Q-values below therefore do not measure the values of the moves that were played.

\begin{table}[htbp]
\centering
\caption{Self-play performance comparison between Vanilla DQN with LSTM and Rainbow DQN with AAREN.}
\label{tab:arch_comparison}
\begin{tabular}{l c c}
\hline
\textbf{Metric} & \textbf{DQN+LSTM (15k)} & \textbf{Rainbow+AAREN (9.5k)} \\ \hline
Total Episodes         & 15,302   & 9,778    \\
Total Steps            & 10,306,433 & 5,169,892 \\
Win Rate (\%)          & 24.97    & 49.88    \\
Draw Rate (\%)         & 50.04    & 6.69     \\
Loss Rate (\%)         & 24.99    & 43.43    \\
Steps per Win          & 2,697    & 1,060    \\
Avg Episode Length     & 1,345    & 1,055    \\
Flag Capture Ratio (\%)& 18.79   & 8.88     \\
Loss--Win Rate Corr.   & 0.051   & $-$0.787 \\
Mean Q-Value           & 0.33    & 0.70     \\
Model Size (per ckpt)  & $\sim$39\,MB & $\sim$433\,MB \\
\hline
\end{tabular}
\end{table}

\subsection{Win Rate }\label{finding:7}
The Rainbow agent won 49.88\% of its self-play games at 9,778 episodes, against 24.97\% for the DQN at 15,302 episodes. In self-play this number is set mostly by the draw rate. When both sides play the same policy, the decisive games split about evenly between them, as they did in both runs: the DQN won 24.97\% and lost 24.99\%, and the Rainbow agent won 49.88\% and lost 43.43\%. The win rate doubled because the share of decisive games nearly doubled (from 50\% to 93\%), not because the agent became stronger than its opponent. Showing that the Rainbow agent plays better than the DQN would need games between the two, or against the same fixed opponents.

\subsection{Draw Elimination}\label{finding:8}
The normal DQN drew 50.04\% of its games, finishing the 1500 moves per agent in an episode. The Rainbow agent drew 6.69\% over 1500 moves. It lost more games outright (43.43\% versus 24.99\%), but it produced definitive outcomes. An agent that draws half of its games has learned to avoid playing passively and conservative in calculating outcomes in the short term. Section~\ref{finding:6} (Draw Dominance) attributed the draws to scalar value collapse. These runs cannot tell which change removed them, because the configurations differ in more than the value head. The effect also did not carry over to the 75k run against heuristic opponents, which drew 77.02\% of its games (Table~\ref{tab:rdqn75k}).

A 43.43\% loss rate alongside a 6.69\% draw rate could suggest that the Rainbow agent takes excessive risks. In self-play, however, every decisive game has one winner and one loser, and both are copies of the same agent. The 93.31\% decisive outcomes mean that the policy forces resolutions rather than stalling. They say nothing about whether its play is reckless or sound against an opponent that is not itself. The self-play games were also shorter (1,055 moves on average, against 1,345 for the normal DQN).

\subsection{Loss--Performance Correlation}\label{finding:9}
The correlation between windowed policy loss and win rate was 0.051 for the normal DQN and $-$0.787 for the Rainbow agent. Two curves that both drift with training time correlate largely through that common drift, whatever their relation. The later formal\_4 run (Rainbow DQN, AAREN off) shows how far this goes. Over blocks of 100 episodes its loss and training win rate correlated at $-0.49$, but the correlation fell to $-0.14$ once a linear trend was removed from both. Its loss fell from 3.63 to 3.05 while its greedy evaluations against fixed opponents did not improve (Appendix~\ref{app:lessons}). The $-0.787$ was not computed on detrended curves, and in self-play the win rate mostly reflects the draw rate (Section~\ref{finding:7}). The coefficient therefore does not show that the Rainbow loss tracks the quality of play. The decoupling of Section~\ref{finding:2} appears in the Rainbow learner as well.

\subsection{Sample Efficiency}\label{finding:10}
The Rainbow agent required 1,060 environment steps per win compared to 2,697 for the normal DQN. At the episode level, the Rainbow agent won once every 2.0 episodes while the normal DQN won once every 4.0, meaning the DQN architecture stumbled its way to a win due to the maximum step is 3000 steps per episode. The Rainbow model consumed 5.17 million total steps; the normal DQN consumed 10.31 million. Steps per win is the total number of steps divided by the number of wins, and in self-play the win rate rose with the share of decisive games (Section~\ref{finding:7}). The metric therefore reflects fewer draws and shorter games more than it reflects learning more from each sample.

\subsection{AAREN and LSTM Accuracy}\label{finding:11}
The AAREN module achieved 24.23\% piece identity inference accuracy, compared to 17.74\% for the LSTM-based PBS module. Three things limit what this gap shows.

First, the 8.3\% baseline of Section~\ref{finding:4} is a uniform guess over 12 types, but the types are not equally common. Scouts are 8 of the 40 pieces (20\%) and 8 of the 33 movable ones (24\%). A predictor that always names the most common type would therefore score about 20 to 24\%, depending on which pieces battles reveal. Until both accuracies are compared with that majority baseline on the same reveals, 24.23\% is not yet evidence of inference from behaviour.

Second, the two accuracies come from different runs, with different opponents, lengths and networks. Third, in the thesis code AAREN read the padding step of every history shorter than the longest in its batch instead of its last real step, and the DQN loss did not train it (Chapter~\ref{chap:methodology}, Gradient Bridge). Section~\ref{finding:4} (LSTM Accuracy Ceiling) predicted that attention would help, but this result neither confirms nor refutes that prediction.

The Q-values cannot be compared either. The LSTM Q-values settled near 0.18 and the Rainbow ones reached 0.84, but the Rainbow trainer's average Q was a maximum over mostly untrained actions.

\subsection{Learned Exploration Active}\label{finding:12}
The NoisyNet entropy rose steadily from 0.034 to 0.049 throughout Rainbow training, adapting to the need for exploration during the training as the agent have not converge to a solution. The normal DQN has a fixed entropy of zero (Section~\ref{finding:5}, Exploration Collapse), due to the constant decay of $\epsilon$-greedy. Noisy Networks replaced frozen $\epsilon$-greedy decay with adaptive, state-dependent exploration.

\subsection{Extended Training at 75,000 Episodes}

A separate Rainbow DQN with AAREN training run reached 75,047 episodes against heuristic opponents. This run used a different curriculum configuration and did not advance to self-play, so it cannot be directly compared to the self-play results above. It does, however, validate two architectural properties over a longer horizon. Table~\ref{tab:rdqn75k} summarizes its metrics.

\begin{table}[htbp]
\centering
\caption{Rainbow DQN with AAREN extended training (75k episodes) against heuristic opponents.}
\label{tab:rdqn75k}
\begin{tabular}{l c}
\hline
\textbf{Metric} & \textbf{Value} \\ \hline
Total Episodes          & 75,047 \\
Win Rate                & 11.66\% \\
Draw Rate               & 77.02\% \\
vs Smart Heuristic WR   & 11.74\% (38,462 games) \\
vs Frozen Heuristic WR  & 11.59\% (43,331 games) \\
Policy Loss (late)      & 2.84 \\
Q-Value (late)          & 1.84 \\
AAREN Accuracy          & 20.78\% \\
Flag Capture Ratio      & 99.96\% \\
Avg Episode Length      & 378 \\
\hline
\end{tabular}
\end{table}

\subsubsection{Loss Convergence Over Long Horizons}\label{finding:13}
The distributional loss fell from 3.89 to 2.84 across 75,000 episodes. The normal DQN showed no comparable long-term loss reduction; its loss either plateaued or inflated in late training. A C51 loss of $\ln 51 \approx 3.93$ is what a prediction spread evenly over the 51 atoms scores, so the run started from an almost uniform prediction. The falling loss shows the critic fitting its own targets more closely. A gradual increase in win rate was seen across the second half of the run, but the run's win rate against heuristic opponents was 11.66\% overall. As in formal\_4 (Section~\ref{finding:9}), the loss can fall while the agent keeps losing, so the falling loss alone is not evidence that extended training improves play.

\subsubsection{Flag Pursuit Behavior}\label{finding:14}
The average episode length in the 75k run settled near 378 steps, compared to 1,345 for the normal DQN and 1,055 for the Rainbow self-play run. The shorter episodes indicate that the agent ended games faster. Of its wins, 99.96\% came by capturing the opponent flag rather than by attrition. The normal DQN won by flag capture only 18.79\% of the time, relying instead on piece depletion. The logged average Q-value grew from near zero to 1.84. With a win worth $+1$, a value of 1.84 at an 11.66\% win rate cannot be the value of the moves played; it is the maximum over mostly untrained actions described above. The attention layer may help the agent locate the flag, but no run has isolated its contribution.

The 75k run is also the more informative measure of strength. In self-play a win rate near 50\% follows from symmetry (Section~\ref{finding:7}). Against fixed heuristic opponents, the agent won about one game in nine (11.66\%) and drew 77.02\% across 75,047 episodes. It did as well against the smart heuristic (11.74\%) as against the frozen one (11.59\%). The later formal runs (Appendix~\ref{app:lessons}) also struggled against these opponents. formal\_4 still lost as many games as it won against a random opponent at episode 3,000, mostly because it attacked revealed pieces it could not beat.

\section{Evidence for Architectural Necessity}

Each component in the thesis title maps to a specific failure mode observed in the normal DQN baseline. Table~\ref{tab:title_evidence} summarizes this mapping and the evidence so far from the Rainbow DQN runs.

\begin{table}[htbp]
    \centering
    \caption{MARQ title components mapped to normal DQN failure modes and the Rainbow DQN evidence so far.}
    \label{tab:title_evidence}
    \begin{tabular}{|p{2.8cm}|p{4.5cm}|p{6.5cm}|}
        \hline
        \textbf{Title Component} & \textbf{Observed Failure Mode} & \textbf{Empirical Evidence} \\
        \hline
        Multi-Agent & Symmetric MARL Nash lock-in & P1 and P2 win rates were 25.1\% and 24.9\% in the normal DQN. In Rainbow self-play they were 49.88\% and 43.43\%, still close, as self-play implies; what changed is that draws fell from 50\% to 6.69\% (Section~\ref{finding:7}). \\
        \hline
        Rainbow DQN & Q-value overestimation, scalar value collapse, exploration death, sparse reward propagation & Self-play draw rate dropped from 50\% to 6.69\% (Section~\ref{finding:8}), but the 75k run drew 77\%. NoisyNet entropy rose throughout training (Section~\ref{finding:12}). The loss--win correlation ($r = -0.787$) was not detrended and does not show coupling (Section~\ref{finding:9}). \\
        \hline
        AAREN & Sequential information bottleneck & AAREN accuracy 24.23\% versus LSTM's 17.74\% (Section~\ref{finding:11}), not yet compared with a majority-class baseline of about 20--24\%. \\
        \hline
        Imperfect Information & Risk-averse draw-seeking policy & The self-play draw rate fell from 38.5\% (Section~\ref{finding:6}) to 6.69\% (Section~\ref{finding:8}). Against heuristic opponents the 75k run drew 77.02\% and won 11.66\%, 99.96\% of its wins by flag capture (Section~\ref{finding:14}). \\
        \hline
    \end{tabular}
\end{table}

The combined findings are consistent with the MARQ architecture, but they do not yet establish what each component contributes. Each component targets a distinct bottleneck. Draw dominance (Section~\ref{finding:6}) is plausibly a compound effect of scalar value collapse, exploration death and sparse reward propagation. Under the combined changes the self-play draw rate fell from 50.04\% to 6.69\% (Section~\ref{finding:8}). No ablation separates the effects of those changes, however, and the run against heuristic opponents still drew 77\% of its games.


\section{Supporting Evidence from Related Work}

These limitations are consistent with findings elsewhere in deep RL for imperfect information games.

P\'{e}rolat et al. (2022) solved full-size Stratego with DeepNash \cite{perolat2022stratego} using game-theoretic regularization (R-NaD), not Q-value maximization. DeepNash shows Stratego is solvable by deep RL, but only when the architecture explicitly handles non-stationarity and strategic diversity. The Ataraxos system \cite{Ataraxos2025} extended this to no-press Diplomacy with KL-divergence stabilization, a mechanism similar to MARQ's dynamic damping but absent from the normal DQN baseline.

Hessel et al. (2021) showed in the Rainbow DQN study that Double DQN, PER, Dueling Networks, Noisy Nets, C51, and multi-step returns each contribute additive gains \cite{RainbowJMLR2021}. Removing any single component degraded performance. No one fix is enough.


\section{Implications for MARQ}

The analysis attributes the normal DQN's performance ceiling to identifiable architectural properties. Table~\ref{tab:limitation_mapping} maps each finding to its theoretical cause, the MARQ component meant to address it, and how far the data so far supports that.

\begin{table}[htbp]
    \centering
    \caption{Observed limitations mapped to theoretical causes, MARQ components, and validation status.}
    \label{tab:limitation_mapping}
    \begin{tabular}{|p{3cm}|p{3.5cm}|p{3.5cm}|p{2.5cm}|}
        \hline
        \textbf{Finding} & \textbf{Theoretical Cause} & \textbf{MARQ Component} & \textbf{Status} \\
        \hline
        Q-Value Overestimation & Maximization Bias & Double DQN / Dueling & Not tested \\
        \hline
        Loss--Perf. Decoupling & Scalar Value Collapse & C51 Distributional RL & Not supported (Sec.~\ref{finding:9}) \\
        \hline
        LSTM Accuracy Ceiling & Sequential Bottleneck & AAREN Attention & Inconclusive (Sec.~\ref{finding:11}) \\
        \hline
        Exploration Collapse & Fixed $\epsilon$-greedy & Noisy Networks & Supported (Sec.~\ref{finding:12}) \\
        \hline
        Reward Propagation & 1-step Bootstrapping & $n$-step Returns & Not tested \\
        \hline
        Draw Dominance & Compound Effect & Full MARQ Pipeline & Self-play only (Sec.~\ref{finding:8}) \\
        \hline
    \end{tabular}
\end{table}

``Not tested'' means that no run isolated the component. Overestimation would need the values of the moves played compared with the returns that followed. The trainer's logged Q was a maximum over mostly untrained actions, and the win rate of Section~\ref{finding:7} does not measure bias. The steps-per-win figure of Section~\ref{finding:10} reflects the draw rate, not how far rewards propagate.

The normal DQN baseline in this chapter served as a reference point. In self-play the Rainbow agent with AAREN produced far fewer draws (6.69\% against 50.04\%) and shorter games, and its exploration noise stayed active. Against heuristic opponents it won 11.66\% of 75,047 games. Two questions remain open. What each component contributes needs ablations evaluated against fixed opponents. Whether AAREN supplies information the observation does not needs runs with and without it, compared on the same evaluations. The formal runs of Appendix~\ref{app:lessons} begin that work, so far with AAREN switched off.
